\documentclass{article}

\usepackage{booktabs}
\usepackage{tabularx}
\usepackage{hyperref}
\usepackage{pdflscape}
\usepackage{longtable}
\usepackage{array}

\hypersetup{
    colorlinks=true,       % false: boxed links; true: colored links
    linkcolor=red,          % color of internal links (change box color with linkbordercolor)
    citecolor=green,        % color of links to bibliography
    filecolor=magenta,      % color of file links
    urlcolor=cyan           % color of external links
}

\title{Hazard Analysis\\\progname}

\author{\authname}

\date{}

\input{../Comments.text}
\input{../Common.text}

% --- REQUIREMENT COUNTERS ---
\newcounter{ENV}
\newcounter{SFR}
\newcounter{SNFR}
\newcounter{SSR}
\renewcommand{\theENV}{ENV-\reqnumber{\value{ENV}}}
\renewcommand{\theSFR}{SFR-\reqnumber{\value{SFR}}}
\renewcommand{\theSNFR}{SNFR-\reqnumber{\value{SNFR}}}
\renewcommand{\theSSR}{SSR-\reqnumber{\value{SSR}}}

% Format numbers as 001, 002, 003, etc.
\newcommand{\reqnumber}[1]{%
	\ifnum#1<10 00%
	\else\ifnum#1<100 0%
	\fi\fi
	\number#1%
}

% --- REQUIREMENT MACRO SETUP ---
\newcommand{\requirement}[6]{%
	\par\noindent
	\begingroup
	
	% Increment the appropriate counter
	\refstepcounter{#1}\label{#6}%
	\edef\reqid{\csname the#1\endcsname}%
	
	
	% Adjustable formatting
	\def\reqindent{2.3cm}
	\setlength{\parindent}{0pt}
	
	% Requirement ID and description
	\noindent
	\makebox[0pt][l]{\textbf{\reqid:}}%
	\hspace*{\reqindent}%
	\parbox[t]{\dimexpr\linewidth-\reqindent\relax}{%
		\parindent=0pt #2%
	}%
	\par\vspace{8pt}
	
	% Indented fields
	\leftskip=\reqindent
	\noindent\textbf{Rationale:} #3\par\vspace{5pt}
	\noindent\textbf{Fit Criterion:} #4\par\vspace{5pt}
	\noindent\textbf{Traceability:} #5\par
	
	\endgroup
	\vspace{0.8em}
}
% --- HOW TO USE IT ---

%\requirement{Tag/ID}
%{Description}
%{Rationale}
%{Fit Criterion}
%{Traceability tags}
%{Unique reference name}

\begin{document}

\maketitle
\thispagestyle{empty}

~\newpage

\pagenumbering{roman}

\begin{table}[hp]
\caption{Revision History} \label{TblRevisionHistory}
\begin{tabularx}{\textwidth}{llX}
\toprule
\textbf{Date} & \textbf{Developer(s)} & \textbf{Change}\\
\midrule
October 17, 2026 & Andreas Sideris & AI Disclosure Updated\\
October 18, 2026 & Andreas Sideris & FMEA Table Completed\\
October 18, 2026 & Andreas Sideris & Reflection Added\\
\bottomrule
\end{tabularx}
\end{table}

~\newpage

\tableofcontents

\listoftables

~\newpage

\pagenumbering{arabic}

\section{Introduction}

\wss{You can include your definition of what a hazard is here.}

\section{Scope and Purpose of Hazard Analysis}

\wss{You should say what \textbf{loss} could be incurred because of the
hazards.}

\section{System Boundaries and Components}

\wss{Dividing the system into components will help you brainstorm the hazards.
You shouldn't do a full design of the components, just get a feel for the major
ones.  For projects that involve hardware, the components will typically include
each individual piece of hardware.  If your software will have a database, or an
important library, these are also potential components.}

\section{Critical Assumptions}

\wss{These assumptions that are made about the software or system.  You should
minimize the number of assumptions that remove potential hazards.  For instance,
you could assume a part will never fail, but it is generally better to include
this potential failure mode.}

\newpage

\begin{landscape}

\section{Failure Mode and Effect Analysis}

\input{SectionFMEA.text}

\end{landscape}

\section{Safety and Security Requirements}

\input{SafetyAndSecurityReqs.text}

\section{Roadmap}

\wss{Which safety requirements will be implemented as part of the capstone timeline?
Which requirements will be implemented in the future?}

\newpage{}

\section{Appendix --- Generative AI Use Disclosure}

\subsection*{AI Tools Used}

\begin{itemize}
	\item OpenAI GPT-6
	\item Google Search AI Overview (Google Gemini)
\end{itemize}

\subsection*{How and Where Used}

AI tools were used for assistance with LaTeX syntax and commands, during 
the development of the FMEA table section, and when developing 
LaTeX formatting for the document as a whole. Specifically, AI tools were 
used for assistance in creating the LaTeX table used for the FMEA summary and 
creating the list of tables, as well as resolving various minor formatting 
bugs. AI tools were also used to create the LaTeX custom requirement format 
macro used in the document, that was also used in the SRS. In this section 
(FMEA), 
AI tools were not used to generate any text or wholly generate any of the 
ideas or conclusions presented.

\subsection*{Verification of AI-Generated Content}

Any LaTeX syntax or code provided by AI tools was verified by testing 
compilation in LaTeX documents, and ensuring the resulting output was 
correctly formatted and matched the needs of the document. No ideas or text 
were newly generated by AI tools and presented in the document.



\newpage{}

\section{Appendix --- Reflection}

\textbf{Reflection Questions Answered:}

\begin{enumerate}
    \item What went well while writing this deliverable? 
    \item What pain points did you experience during this deliverable, and how
    did you resolve them?
    \item Which of your listed risks had your team thought of before this
    deliverable, and which did you think of while doing this deliverable? For
    the latter ones (ones you thought of while doing the Hazard Analysis), how
    did they come about?
    \item Other than the risk of physical harm (some projects may not have any
    appreciable risks of this form), list at least 2 other types of risk in
    software products. Why are they important to consider?
\end{enumerate}

\textbf{Andreas:}

\begin{enumerate}
	\item When writing this deliverable, our team worked well to synchronize 
	our decisions between sections. While I primarily focused on the FMEA 
	table, I was able to discuss decisions my teammates had made on other vital 
	sections (such as System Boundaries and Components) with them, clarify 
	uncertainties, and use the output of our discussions to guide my work.
	\item One pain point I experienced during this deliverable was finding 
	suitable failure modes and hazards for our software, since it is intended 
	to be a video game (it does not have many obvious physical safety risks). 
	This was resolved by considering other ways failure of the system could 
	result in harm, such as it being inaccessible to users with photosensitive 
	epilepsy, or colourblind users. It was also found that other failures of 
	the system could result in harm (in the expanded sense used in this 
	project), by frustrating users and degrading the user experience overall.
	\item Some risks had been considered before, such as failures of the system 
	resulting in corrupted user save data, or unexpected behavior (that could 
	frustrate the user), but specific risks concerning colourblind users and 
	users with photosensitive epilepsy were newly found and formalized. These 
	risks were found by considering the ways in which failures of the system 
	could cause harm to specific user groups, if they were not accounted for in 
	the design (as these risks could be missed without careful consideration). 
	This led to finding a physical safety hazard (the potential for epileptic 
	seizures in those with photosensitive epilepsy), as well as a potential 
	risk of preventing colourblind users from playing the game as others can.
	\item One risk can involve the risk of preventing certain users from being 
	able to use a software product (as discussed above, in regards to 
	accessibility). This risk is important to consider because it can result in 
	real harm; users with specific conditions can feel alienated when they 
	aren't able to use or enjoy software products that others do. Another risk 
	involves the risk of users becoming confused or frustrated when using the 
	software product (as discussed in the FMEA analysis, when users' saved 
	progress is lost, or the user interface is difficult to understand). This 
	must be considered, as if most users become frustrated when using a 
	software product, it will fail commercially/publicly in favour of other 
	products that are designed with the user in mind.
\end{enumerate}

\textbf{Humna:}

\begin{enumerate}
	\item 
	\item 
	\item 
	\item 
\end{enumerate}

\textbf{Christopher:}

\begin{enumerate}
	\item 
	\item 
	\item 
	\item 
\end{enumerate}

\textbf{Andy:}

\begin{enumerate}
	\item 
	\item 
	\item 
	\item 
\end{enumerate}

\textbf{Owen:}

\begin{enumerate}
	\item 
	\item 
	\item 
	\item 
\end{enumerate}

\end{document}